\chapter{Introduction}

\section{Context and Motivation}
The exact string matching problem is one of the fundamental tasks in classical algorithmics and text processing. Given a text $T$ of length $n$ and a pattern $P$ of length $m$ defined over a finite alphabet $\Sigma$, the objective is to locate all valid shifts $i$ such that every character of the pattern aligns perfectly with the corresponding substring of the text. Standard solutions to this problem, such as the Knuth-Morris-Pratt \cite{kmp1977} and Boyer-Moore \cite{bm1977} algorithms, achieve optimal linear time bounds of $\mathcal{O}(n + m)$ through efficient pattern preprocessing and shifting heuristics. Another notable probabilistic approach is the Karp-Rabin algorithm \cite{karp1987efficient}, which utilizes rolling hashes for highly efficient average-case matching. See also \cite{aho1990} for an overview summary of string matching techniques.

However, standard matching models assume that all input symbols are explicitly known. The \textit{string matching with wildcards} problem extends this classic formulation by introducing a special wildcard symbol (commonly denoted as \texttt{?} or ``don't care''). This symbol is strictly defined to successfully match any arbitrary character from the underlying alphabet, as well as itself. Resolving this generalized problem is highly relevant in practical applications, such as bioinformatics and query parsing, which frequently process incomplete data, uncertainty, or intentional omissions \cite{mp1994}.

While conceptually simple, introducing wildcards fundamentally alters the mathematical structure of the matching process. Classical linear-time algorithms rely inherently on the transitive property of character equality. Under wildcard matching rules, this transitivity is explicitly violated:
\begin{equation}
    (a = \text{\texttt{?}}) \land (\text{\texttt{?}} = b) \centernot\implies (a = b) \quad \text{for } a \neq b.
\end{equation}

Because this equivalence relation fails, the deterministic prefix tables and shifting mechanisms utilized by algorithms like Knuth-Morris-Pratt become unusable. Solving the wildcard matching problem via naive character-by-character comparison forces the evaluation of every possible shift individually, resulting in a degraded worst-case time complexity of $\mathcal{O}(n \cdot m)$. For large-scale inputs, such as DNA sequences where $n \sim 10^7$ and $m \sim 10^4$, this quadratic boundary causes a severe computational bottleneck, necessitating the search for alternative approaches \cite{fp1974}.

\section{Historical Evolution}
Overcoming this quadratic barrier required moving away from standard character comparisons and using algebraic convolution techniques.

The first major breakthrough was published by Fischer and Paterson in 1974 \cite{fp1974}. They showed that string matching with wildcards can be translated into boolean vector convolutions. By encoding letters as indicator polynomials, the matching problem turns into polynomial multiplication. When implemented using large integers or discrete transforms, their method achieved an upper bound of $\mathcal{O}(|\Sigma|^2 \cdot n \log n)$. 

While this is asymptotically faster than $\mathcal{O}(n \cdot m)$ for small alphabets, Fischer and Paterson's method has severe practical limits. In practice, multiplying such massive numbers introduces heavy computational overhead. Furthermore, the performance depends heavily on the alphabet size $|\Sigma|$, making this method completely unusable for extended character sets like Unicode.

To fix the alphabet bottleneck, researchers explored two main paths. The first focused on combinatorics. Muthukrishnan and Palem \cite{mp1994} introduced an optimization based on Sperner's theorem, reducing the number of required convolutions from $|\Sigma|$ to $\mathcal{O}(\log |\Sigma|)$. The second path used randomized projections. Indyk \cite{indyk1998} proposed a randomized approach that runs in $\mathcal{O}(n \log n)$ time, completely independent of the alphabet size. Shortly after, Kalai \cite{kalai2002} developed a much simpler algebraic fingerprinting method that achieved the optimal $\mathcal{O}(n \log m)$ complexity bound using only two Fast Fourier Transform (FFT) passes. Finally, Clifford and Clifford \cite{clifford2007} successfully resolved the long-standing quest for an optimal, alphabet-independent deterministic algorithm, achieving the same $\mathcal{O}(n \log m)$ bound without relying on probabilistic error margins.

\section{Research Objectives}
While the asymptotic complexities of these algorithms are well documented in theoretical papers, big-$\mathcal{O}$ notation hides important practical details. It ignores constant factors, memory layout penalties, and the real-world overhead introduced by floating-point arithmetic and hardware limitations. Although these algorithms are thoroughly analyzed in theoretical computer science, detailed benchmarks comparing them side-by-side are less common.

The main goal of this thesis is to evaluate theoretical algorithms in the context of practical software engineering constraints. To achieve this, several key approaches were systematically implemented and compared in Python. The analysis covers a full spectrum of methods, starting from the naive baseline and the classical Fischer-Paterson bitwise method, through FFT-based convolutions (Baseline FFT and Sperner), probabilistic techniques proposed by Indyk and Kalai, and concluding with the deterministic Clifford-Clifford algorithm.

During the implementation, I also had to resolve several practical engineering problems, such as handling vector padding for FFT, managing floating-point inaccuracies, and controlling random error bounds. Finally, detailed performance tests were conducted to see how these algorithms scale with different text lengths, alphabet sizes, and wildcard densities. The results of these tests serve as the basis for a practical comparative framework for future implementations.